\documentclass[a4paper, 12pt]{article}

%%% Работа с русским языком
\usepackage{cmap}					% поиск в PDF
\usepackage{mathtext} 				% русские буквы в формулах
\usepackage[T2A]{fontenc}			% кодировка
\usepackage[utf8]{inputenc}			% кодировка исходного текста
\usepackage[russian]{babel}	% локализация и переносы

%%% Дополнительная работа с математикой
\usepackage{amsmath,amsfonts,amssymb,amsthm,mathtools} % AMS
\usepackage{icomma} % "Умная" запятая: $0,2$ --- число, $0, 2$ --- перечисление

%% Номера формул
%\mathtoolsset{showonlyrefs=true} % Показывать номера только у тех формул, на которые есть \eqref{} в тексте.
%\usepackage{leqno} % Немуреация формул слева

%% Шрифты
\usepackage{euscript}	 % Шрифт Евклид
\usepackage{mathrsfs} % Красивый матшрифт

%%% Свои команды
\DeclareMathOperator{\sgn}{\mathop{sgn}}

%% Поля
\usepackage[left=2cm,right=2cm,top=2cm,bottom=2cm,bindingoffset=0cm]{geometry}

%% Русские списки
\usepackage{enumitem}
\makeatletter
\AddEnumerateCounter{\asbuk}{\russian@alph}{щ}
\makeatother

%%% Работа с картинками
\usepackage{graphicx}  % Для вставки рисунков
\graphicspath{{images/}{images2/}}  % папки с картинками
\setlength\fboxsep{3pt} % Отступ рамки \fbox{} от рисунка
\setlength\fboxrule{1pt} % Толщина линий рамки \fbox{}
\usepackage{wrapfig} % Обтекание рисунков и таблиц текстом

%%% Работа с таблицами
\usepackage{array,tabularx,tabulary,booktabs} % Дополнительная работа с таблицами
\usepackage{longtable}  % Длинные таблицы
\usepackage{multirow} % Слияние строк в таблице

%% Красная строка
\setlength{\parindent}{2em}

%% Интервалы
\linespread{1}
\usepackage{multirow}

%% TikZ
\usepackage{tikz}
\usetikzlibrary{graphs,graphs.standard}

%% Верхний колонтитул
\usepackage{fancyhdr}
\pagestyle{fancy}

%% Перенос знаков в формулах (по Львовскому)
\newcommand*{\hm}[1]{#1\nobreak\discretionary{}
	{\hbox{$\mathsurround=0pt #1$}}{}}

%% дополнения
\usepackage{float} % Добавляет возможность работы с командой [H] которая улучшает расположение на странице
\usepackage{gensymb} % Красивые градусы
\usepackage{caption} % Пакет для подписей к рисункам, в частности, для работы caption*

% подключаем hyperref (для ссылок внутри  pdf)
\usepackage[unicode, pdftex]{hyperref}

%%% Теоремы
\theoremstyle{plain}                    % Это стиль по умолчанию, его можно не переопределять.
\renewcommand\qedsymbol{$\blacksquare$} % переопределение символа завершения доказательства

\newtheorem{theorem}{Теорема}[section] % Теорема (счетчик по секиям)
\newtheorem{proposition}{Утверждение}[section] % Утверждение (счетчик по секиям)
\newtheorem{definition}{Определение}[section] % Определение (счетчик по секиям)
\newtheorem{corollary}{Следствие}[theorem] % Следстиве (счетчик по теоремам)
\newtheorem{problem}{Задача}[section] % Задача (счетчик по секиям)
\newtheorem*{remark}{Примечание} % Примечание (можно переопределить, как Замечание)
\newtheorem{lemma}{Лемма}[section] % Лемма (счетчик по секиям)

\newtheorem{example}{Пример}[section] % Пример
\newtheorem{counterexample}{Контрпример}[section] % Контрпример

\begin{document}
    \newcommand{\HRule}{\rule{\linewidth}{0.7mm}} % Defines a new command for the horizontal lines, change thickness here
	
	\begin{center}
		\large\textbf{Московский Физико-Технический Институт}\\ % Name of your university/college
		\large\textbf{(государственный университет)}
	
		\vfill
		
		\Large Научный проект по защите информации\\[0.5cm] % Preambule of your document title
		
		
		\HRule
		\\[0.4cm]
		{ \huge \bfseries *Реализация шифра <<Кузнечик>> с применением векторных инструкций RISC-V архитектуры*} % Title of your document
		\\[0.4cm]
		\HRule
		\\[0.5cm]
		
		\ \\
	\textbf{\large Авторы:} \\
	\large Филиппенко Павел Б01-009\\
        \large Фролов Даниил Б01-009\\
        \large Долгодворов Егор Б01-009\\
        \large Белов Владислав Б01-009\\
		\vfill
		\hspace*{-0.8 cm}\includegraphics[width=100 pt]{../images/frkt_logo}\\ % Logo of your  company/university/college
		\large Долгопрудный, 2023 % Location and year
	\end{center}

\newpage
\setcounter{page}{2}
\fancyfoot[c]{\thepage}
\fancyhead[L] {Криптография} % Information in page header
\fancyhead[R]{}

    \section{Abstract}
    Для современных компьютерных систем важно, чтобы алгоритмы шифрования данных были не
    только надежными, но так же быстрыми с точки зрения вычислительных операций. В
    данной статье мы опишем векторную реализацию блочного шифра <<Кузнечик>> для 
    процессорной архитектуры RISC-V, а так же сравним скорость ее работы с 
    последовательной реализацией и векторной реализацией на других архитектурах.

    \section{Введение}

    В современном мире, где довольно значительную роль играют цифровые технологии,
    криптография является довольно важной частью обеспечения безопасности в цифровой
    среде. Защита интернет-соединений и каналов связи, надежность проведения банковских
    операций, подтверждение подлинности электронных документов. Современные методы
    защиты информации накладывают большое количество требований на современные 
    алгоритмы шифрования, такие как теоретическая и вычислительная криптостойкость, а 
    так же скорость работы.
    
    Для современных компьютерных систем важно, чтобы алгоритмы шифрования данных были не
    только надежными, но так же быстрыми с точки зрения вычислительных операций. На 
    сегодняшний день одним из популярных способов увеличения скорости работы программ
    является векторизация математических операций. Данный способ основывается на 
    принципе компьютерных вычислений SIMD (анг. single instruction multiple data), 
    позволяющем обеспечить параллелизм на уровне данных. И хотя, данный вид
    оптимизации требует индивидуальной реализации для разных архитектур, он является
    довольно мощным в плане увеличения производительности кода. Кроме того, векторизацию
    удобно использовать в алгоритмах с большим количеством однотипных математических
    операций, в частности, в алгоритмах блочного шифрования.

    В данной статье мы опишем векторную реализацию блочного шифра <<Кузнечик>> для 
    процессорной архитектуры RISC-V, а так же сравним скорость ее работы с 
    последовательной реализацией и векторной реализацией на других архитектурах.

    \section{Шифр <<Кузнечик>>}

    Алгоритм шифрования <<Кузнечик>> является одним из двух шифров, принятых в 2015 г.
    в России в стандарт блочного шифрования ГОСТ Р 34.12-2015. <<Кузнечик>> относится
    к типу блочных симметричных шифров; это значит, что для зашифрования и
    расшифрования используется один и тот же ключ, и что при шифровании исходный текст
    разделяется на блоки фиксированной длины, каждый из которых шифруется отдельно.

    На сегодняшний день, шифр все еще считается криптостойким.

    \section{Математическое описания шифра <<Кузнечик>>}

    

    % \section{Режимы работы.}

\end{document}